\documentclass[11pt]{article}
\usepackage[margin=1in]{geometry}
\usepackage{times}
\usepackage{hyperref}
\hypersetup{colorlinks=true, linkcolor=black, urlcolor=blue, citecolor=black}
\title{A Radio Receiver Is Not a Generator: The Hendershot Fuelless Generator}
\author{Flyxion \\ \small Independent Researcher}
\date{}
\begin{document}
\maketitle

\begin{abstract}
Lester Hendershot demonstrated a device in 1928, publicized in connection with his acquaintance Charles Lindbergh, claimed to generate usable electrical power from Earth's ambient magnetic field using a specially wound coil arrangement requiring no external fuel or battery, a claim that attracted brief national press attention before independent examiners found the device was, in its most notable public tests, receiving power from a nearby radio broadcast antenna and, in a separate later demonstration, from a concealed battery. This paper treats the case as an early and comparatively well-documented instance of the free-energy demonstration genre, notable for the specific, identified conventional power sources found in two independent examinations of the same claim decades apart, and considers why the claim nonetheless persisted in fringe literature well after both explanations were a matter of public record.
\end{abstract}

\section{Introduction}

Lester Hendershot claimed in 1928 to have built a device, a specially configured coil and magnet arrangement, capable of generating usable electrical current directly from Earth's magnetic field with no battery or external power connection, and gained national attention when Charles Lindbergh, a personal acquaintance, expressed public interest in the claim, leading to press coverage and a demonstration arranged before technical examiners at the request of United States Navy and other officials.

\section{The First Examination}

Contemporary technical examination of Hendershot's original demonstration device found that the device's apparent output correlated closely with, and was ultimately traced to, a nearby powerful radio broadcast transmission being picked up by the device's coil acting, in effect, as a crude radio receiver antenna rather than as a generator producing power from Earth's magnetic field independent of any external transmission source, a mundane finding that received considerably less press attention than the initial announcement had, a pattern common across the free-energy demonstration genre, in which the exciting claim travels further and faster than its subsequent quiet correction.

\section{The Later Demonstration and Concealed Battery}

Hendershot continued to promote refined versions of the device in subsequent years, and a later demonstration, examined independently, was found to include a small battery concealed within the device housing, sufficient to power the modest demonstrated output for the duration of the exhibition without requiring any external or magnetic-field energy source at all, a second and independently identified conventional explanation distinct from the radio-reception finding of the original 1928 examination, indicating that at least two different mundane power sources were identified across the device's promotional history at different points, a detail that considerably undercuts any suggestion that a single misidentified but otherwise consistent anomalous effect was present throughout.

\section{Why Two Different Mundane Explanations Across the Device's History Matters}

A genuine anomalous physical effect, if real, would be expected to manifest consistently across demonstrations of what is presented as the same underlying technology, and finding two entirely different, unrelated mundane explanations, incidental radio reception in one case and a deliberately concealed battery in another, for what was marketed as a single consistent invention is more parsimoniously explained by the device never having produced genuine excess energy at any point, with different demonstrations relying on whatever conventional power source was locally convenient or necessary to sustain the appearance of function for that particular audience, than by positing that a real effect was present but happened to be independently misdiagnosed by two separate examinations using two entirely different, mutually exclusive mundane explanations.

\section{Why the Claim Persisted in Fringe Literature Anyway}

Hendershot's association with Lindbergh's brief public endorsement provided the claim with a durable, respectable-sounding anecdotal anchor that continued to be cited in later free-energy literature independent of, and often without mention of, either the radio-reception or concealed-battery findings, illustrating a recurring pattern across the devices surveyed in this series: an early, well-publicized association with a credible or famous figure tends to propagate through subsequent retellings considerably more durably than the specific, contemporaneous technical corrections that followed, since the endorsement is a simple and memorable narrative element while the correction requires engaging with examination details that later retellings frequently omit.

\section{The Entropic Gravity Framing}

The Hendershot generator makes no gravitational claim, proposing energy extraction from Earth's magnetic field rather than gravitational coupling, and is included in this series as an early historical comparative case in the free-energy genre, notable for how cleanly two independent, mutually distinct mundane explanations were identified across its promotional history, a pattern that recurs in less cleanly documented form across many later devices surveyed elsewhere in this series.

\section{Objections}

It might be argued that the specific concealed-battery finding applied only to a later demonstration and does not retroactively explain the original 1928 device examined separately and found to be receiving radio transmissions, and that treating both findings as jointly discrediting the underlying invention conflates two distinct devices built years apart. This is a fair distinction to draw between the specific instances, but it does not rehabilitate the underlying claim; if anything, it means that across the device's public history, no single demonstrated version was found, upon independent examination, to produce energy through any mechanism other than a conventional and previously undisclosed source, whether an external radio signal or a hidden battery, at every point that independent examination was actually conducted.

\section{Conclusion}

Two separate, independently conducted examinations of Hendershot's fuelless generator across different points in its promotional history each identified a distinct, ordinary explanation for the device's apparent output, radio reception in one case and a concealed battery in another, and the claim's persistence in fringe literature well after both findings were public record illustrates how a credible early endorsement can outlive the specific technical corrections that followed it by decades.

\end{document}
